% VI/22 editorial reader bridge — additive, not a canonical problem.
\noindent\textbf{Theory and worked example: open versus closed subschemes.}
For an affine scheme $X=\operatorname{Spec}A$, an open distinguished subscheme
$D(f)$ has coordinate ring $A_f$: the element $f$ becomes invertible. In
contrast, an ideal $I\subseteq A$ defines a closed subscheme
$\operatorname{Spec}(A/I)$. A closed subscheme contains more information than
its underlying closed subset, since the quotient retains nilpotents.

For instance, in $\mathbb A^1_k=\operatorname{Spec}k[t]$ the ideals $(t)$
and $(t^2)$ define the same closed point set, but different schemes:
\[
Z_{\mathrm{red}}=\operatorname{Spec}k[t]/(t)
\cong\operatorname{Spec}k,
\qquad
Z_{\mathrm{thick}}=\operatorname{Spec}k[t]/(t^2).
\]
The element $\overline t$ is nonzero and nilpotent in the second quotient.
Meanwhile, the complementary open subscheme $D(t)$ is
$\operatorname{Spec}k[t,t^{-1}]$: neither closed quotient is the coordinate
ring of this open part. These constructions explain why ideal sheaves,
localization and the scheme structure must be kept separate.

\medskip
